\begin{proposition}
\label{proposition-cohomological-dimension-spectral}
\begin{reference}
Part (1) is the main theorem of \cite{Scheiderer}.
\end{reference}
Let $X$ be a spectral space of Krull dimension $d$.
Let $\mathcal{F}$ be an abelian sheaf on $X$.
\begin{enumerate}
\item $H^q(X, \mathcal{F}) = 0$ for $q > d$,
\item $H^d(X, \mathcal{F}) \to H^d(U, \mathcal{F})$ is surjective
for every quasi-compact open $U \subset X$,
\item $H^q_Z(X, \mathcal{F}) = 0$ for $q > d$ and any constructible
closed subset $Z \subset X$.
\end{enumerate}
\end{proposition}

\begin{proof}
We prove this result by induction on $d$.

\medskip\noindent
If $d = 0$, then $X$ is a profinite space, see
Topology, Lemma \ref{topology-lemma-characterize-profinite-spectral}.
Thus (1) holds by Lemma \ref{lemma-vanishing-for-profinite}.
If $U \subset X$ is quasi-compact open, then $U$ is
also closed as a quasi-compact subset of a Hausdorff space.
Hence $X = U \amalg (X \setminus U)$ as a topological space
and we see that (2) holds. Given $Z$ as in (3) we consider the
long exact sequence
$$
H^{q - 1}(X, \mathcal{F}) \to
H^{q - 1}(X \setminus Z, \mathcal{F}) \to
H^q_Z(X, \mathcal{F}) \to H^q(X, \mathcal{F})
$$
Since $X$ and $U = X \setminus Z$ are profinite (namely $U$ is quasi-compact
because $Z$ is constructible) and since
we have (2) and (1) we obtain the desired vanishing of the
cohomology groups with support in $Z$.

\medskip\noindent
Induction step. Assume $d \geq 1$ and assume
the proposition is valid for all spectral
spaces of dimension $< d$. We first prove part (2) for $X$.
Let $U$ be a quasi-compact open. Let $\xi \in H^d(U, \mathcal{F})$.
Set $Z = X \setminus U$. Let $W \subset X$ be the set of points
specializing to $Z$. By
Lemma \ref{lemma-cohomology-of-neighbourhoods-of-closed} we have
$$
H^d(W \setminus Z, \mathcal{F}|_{W \setminus Z}) =
\colim_{Z \subset V} H^d(V \setminus Z, \mathcal{F})
$$
where the colimit is over the quasi-compact open neighbourhoods $V$
of $Z$ in $X$.
By Topology, Lemma \ref{topology-lemma-make-spectral-space} we see that
$W \setminus Z$ is a spectral space.
Since every point of $W$ specializes to a point of $Z$, we see that
$W \setminus Z$ is a spectral space of Krull dimension $< d$.
By induction hypothesis we see that the image of $\xi$ in
$H^d(W \setminus Z, \mathcal{F}|_{W \setminus Z})$ is zero.
By the displayed formula, there exists a $Z \subset V \subset X$
quasi-compact open such that $\xi|_{V \setminus Z} = 0$.
Since $V \setminus Z = V \cap U$ we conclude by the Mayer-Vietoris
(Lemma \ref{lemma-mayer-vietoris}) for the covering $X = U \cup V$
that there exists a $\tilde \xi \in H^d(X, \mathcal{F})$ which restricts
to $\xi$ on $U$ and to zero on $V$. In other words, part (2) is true.

\medskip\noindent
Proof of part (1) assuming (2). Choose an injective resolution
$\mathcal{F} \to \mathcal{I}^\bullet$. Set
$$
\mathcal{G} = \Im(\mathcal{I}^{d - 1} \to \mathcal{I}^d) =
\Ker(\mathcal{I}^d \to \mathcal{I}^{d + 1})
$$
For $U \subset X$ quasi-compact open we have a map of exact sequences
as follows
$$
\xymatrix{
\mathcal{I}^{d - 1}(X) \ar[r] \ar[d] &
\mathcal{G}(X) \ar[r] \ar[d] &
H^d(X, \mathcal{F}) \ar[d] \ar[r] & 0 \\
\mathcal{I}^{d - 1}(U) \ar[r] &
\mathcal{G}(U) \ar[r] &
H^d(U, \mathcal{F}) \ar[r] & 0
}
$$
The sheaf $\mathcal{I}^{d - 1}$ is flasque by
Lemma \ref{lemma-injective-flasque} and the fact that $d \geq 1$.
By part (2) we see that the right vertical arrow is surjective.
We conclude by a diagram chase that the map
$\mathcal{G}(X) \to \mathcal{G}(U)$ is surjective.
By Lemma \ref{lemma-vanishing-ravi} we conclude that
$\check{H}^q(\mathcal{U}, \mathcal{G}) = 0$ for $q > 0$ and
any finite covering $\mathcal{U} : U = U_1 \cup \ldots \cup U_n$
of a quasi-compact open by quasi-compact opens. Applying
Lemma \ref{lemma-cech-vanish-basis} we find that $H^q(U, \mathcal{G}) = 0$
for all $q > 0$ and all quasi-compact opens $U$ of $X$.
By Leray's acyclicity lemma
(Derived Categories, Lemma \ref{derived-lemma-leray-acyclicity})
we conclude that
$$
H^q(X, \mathcal{F}) =
H^q\left(
\Gamma(X, \mathcal{I}^0) \to \ldots \to
\Gamma(X, \mathcal{I}^{d - 1}) \to \Gamma(X, \mathcal{G})
\right)
$$
In particular the cohomology group vanishes if $q > d$.

\medskip\noindent
Proof of (3).  Given $Z$ as in (3) we consider the long exact sequence
$$
H^{q - 1}(X, \mathcal{F}) \to
H^{q - 1}(X \setminus Z, \mathcal{F}) \to
H^q_Z(X, \mathcal{F}) \to H^q(X, \mathcal{F})
$$
Since $X$ and $U = X \setminus Z$ are spectral spaces
(Topology, Lemma \ref{topology-lemma-spectral-sub})
of dimension $\leq d$
and since we have (2) and (1) we obtain the desired vanishing.
\end{proof}